\chapter{1981年全国卷（文）}

\section{解答题}

\begin{problem}
设 \(A\) 表示有理数的集合，\(B\) 表示无理数的集合，即设 \(A = \{\text{有理数}\}\)，\(B = \{\text{无理数}\}\)，试写出：

\begin{enumerate}
\item \(A \cup B\)；
\item \(A \cap B\).
\end{enumerate}
\end{problem}

\begin{solution}
\begin{enumerate}
\item 任意实数或者是有理数，或者是无理数，因此 \(A\cup B\) 包含全体实数；反过来，\(A\) 和 \(B\) 都是实数的子集，所以 \(A\cup B\) 不包含实数以外的数. 故 \(A\cup B=\R\).
\item 一个数不可能同时是有理数和无理数，因此 \(A\) 与 \(B\) 没有公共元素. 故 \(A\cap B=\varnothing\).
\end{enumerate}
\end{solution}

\begin{problem}
化简： \(\left[-\frac{a^7b^2}{\sqrt{3}(a+b)^2}\right]^2\times \left[\frac{a^2-b^2}{a^2\sqrt{b}}\right]^4\div \left[\frac{a^2(b-a)}{2}\right]^3\).
\end{problem}

\begin{solution}
原式有意义时，需满足 \(b>0\)、\(a\ne0\)、\(a+b\ne0\) 且 \(a\ne b\). 先分别化简各因式，得
\(\left[-\frac{a^7b^2}{\sqrt{3}(a+b)^2}\right]^2=\frac{a^{14}b^4}{3(a+b)^4}\)，\(\left[\frac{a^2-b^2}{a^2\sqrt b}\right]^4=\frac{(a^2-b^2)^4}{a^8b^2}\)
以及
\[
\left[\frac{a^2(b-a)}{2}\right]^3=\frac{a^6(b-a)^3}{8}
\]
所以原式
\[
=\frac{a^{14}b^4}{3(a+b)^4}\cdot\frac{(a^2-b^2)^4}{a^8b^2}\cdot\frac{8}{a^6(b-a)^3}
=\frac{8b^2(a^2-b^2)^4}{3(a+b)^4(b-a)^3}
\]
由 \(a^2-b^2=(a-b)(a+b)=-(b-a)(a+b)\)，有
\[
(a^2-b^2)^4=(b-a)^4(a+b)^4
\]
约去非零因式后，原式 \(=\dfrac{8}{3}(b-a)b^2\).
\end{solution}

\begin{problem}
在 A，B，C，D 四位候选人中，

\begin{enumerate}
\item 如果选举正，副班长各 \(1\text{ 人}\)，共有几种选法？写出所有可能的选举结果；
\item 如果选举班委 \(3\text{ 人}\)，共有几种选法？写出所有可能的选举结果.
\end{enumerate}
\end{problem}

\begin{solution}
\begin{enumerate}
\item 正班长有 \(4\) 种选法；正班长确定后，副班长有 \(3\) 种选法. 因此选法种数为
\[
4\times3=12
\]
以下把两个字母依次记为正班长和副班长，所有可能的选举结果为 AB、AC、AD、BA、BC、BD、CA、CB、CD、DA、DB、DC.
\item 只需从四位候选人中选出三人，先不区分三人的顺序，因此选法种数为
\[
\mathrm C_4^3=\frac{4\times3\times2}{3\times2\times1}=4
\]
所有可能的选举结果为 ABC、ABD、ACD、BCD.
\end{enumerate}
\end{solution}

\begin{problem}
求函数 \(f(x)=\sin x+\cos x\) 在区间 \((-\pi,\pi)\) 上的最大值.
\end{problem}

\begin{solution}
由三角函数的和角公式，
\[
f(x)=\sin x+\cos x=\sqrt{2}\sin\left(x+\frac{\pi}{4}\right)
\]
因为正弦函数的值不超过 \(1\)，所以对区间内任意 \(x\)，都有 \(f(x)\le\sqrt{2}\). 当 \(x=\dfrac{\pi}{4}\) 时，\(x\in(-\pi,\pi)\)，且
\[
f\left(\frac{\pi}{4}\right)=\sin\frac{\pi}{4}+\cos\frac{\pi}{4}=\sqrt{2}
\]
因此，函数在给定区间上的最大值为 \(\sqrt{2}\).
\end{solution}

\begin{problem}
写出正弦定理，并对钝角三角形的情况加以证明.
\end{problem}

\begin{solution}
在 \(\triangle ABC\) 中，记 \(BC=a\)、\(CA=b\)、\(AB=c\)，则正弦定理为
\[
\frac{a}{\sin A}=\frac{b}{\sin B}=\frac{c}{\sin C}
\]

不妨设钝角为 \(C\)，即 \(C>\dfrac{\pi}{2}\). 从点 \(A\) 向直线 \(BC\) 作垂线，垂足为 \(D\)；从点 \(B\) 向直线 \(AC\) 作垂线，垂足为 \(E\). 因为 \(C\) 是钝角，\(D\) 在射线 \(CB\) 的反向延长线上，\(E\) 在射线 \(CA\) 的反向延长线上.

在直角三角形 \(ACD\) 和 \(ABD\) 中，分别有
\(AD=b\sin(\pi-C)=b\sin C\)，\(AD=c\sin B\).
于是
\[
\frac{b}{\sin B}=\frac{c}{\sin C}
\]
又在直角三角形 \(ABE\) 和 \(BCE\) 中，分别有
\(BE=c\sin A\)，\(BE=a\sin(\pi-C)=a\sin C\).
于是
\[
\frac{a}{\sin A}=\frac{c}{\sin C}
\]
合并两式，即得
\[
\frac{a}{\sin A}=\frac{b}{\sin B}=\frac{c}{\sin C}
\]
\end{solution}

\begin{problem}
已知正方形 \(ABCD\) 的相对顶点 \(A(0,-1)\) 和 \(C(2,5)\)，求顶点 \(B\) 和 \(D\) 的坐标.
\end{problem}

\begin{solution}
设正方形两条对角线的交点为 \(O\). 因为对角线互相平分，
\[
O=\left(\frac{0+2}{2},\frac{-1+5}{2}\right)=(1,2)
\]
于是
\[
\overrightarrow{OA}=(-1,-3)
\]
正方形的两条对角线互相垂直且长度相等，所以从 \(O\) 指向 \(B\) 或 \(D\) 的向量应当与 \(\overrightarrow{OA}\) 垂直，且长度相等. 将 \(\overrightarrow{OA}\) 旋转 \(90^{\circ}\)，得到两个可能的向量 \((3,-1)\) 和 \((-3,1)\)，故两个端点的无序集合为 \(\{(4,1),(-2,3)\}\). 按 \(A\)，\(B\)，\(C\)，\(D\) 逆时针编号的通常约定，取 \(\overrightarrow{OB}=(3,-1)\)、\(\overrightarrow{OD}=(-3,1)\)，因而
\(B=O+(3,-1)=(4,1)\)，\(D=O+(-3,1)=(-2,3)\).
故顶点 \(B\) 和 \(D\) 的坐标分别为 \(B(4,1)\) 和 \(D(-2,3)\).
\end{solution}

\begin{problem}
设 \(1980\) 年底我国人口以 \(10\text{ 亿}\) 计算.

\begin{enumerate}
\item 如果我国人口每年比上年平均递增 \(2\%\)，那么到 \(2000\) 年底将达到多少？
\item 要使 \(2000\) 年底我国人口不超过 \(12\text{ 亿}\)，那么每年比上年平均递增率最高是多少？
\end{enumerate}

\begin{center}
\fbox{%
\begin{minipage}{10cm}
\centering
下列对数值可供选用：

{
\small
\setlength{\tabcolsep}{2pt}
\begin{tabular}{@{}ccc@{}}
\(\lg 1.0087 = 0.00377\) & \(\lg 1.0092 = 0.00396\) & \(\lg 1.0096 = 0.00417\) \\
\(\lg 1.0200 = 0.00860\) & \(\lg 1.2000 = 0.07918\) & \(\lg 1.3098 = 0.11720\) \\
\(\lg 1.4568 = 0.16340\) & \(\lg 1.4859 = 0.17200\) & \(\lg 1.5157 = 0.18060\)
\end{tabular}
}
\end{minipage}%
}
\end{center}
\end{problem}

\begin{solution}
\begin{enumerate}
\item 设人口数的单位为亿. 从 \(1980\) 年底到 \(2000\) 年底共有 \(20\) 年，因此人口数为等比数列的第 \(21\) 项：
\[
N=10(1+0.02)^{20}=10(1.0200)^{20}
\]
取常用对数，利用题给数据得
\[
\lg N=1+20\lg1.0200=1+20\times0.00860=1.17200
\]
又因为 \(\lg1.4859=0.17200\)，所以
\[
N=10\times1.4859=14.859
\]
故到 \(2000\) 年底我国人口将达到 \(14.859\text{ 亿}\).
\item 设每年比上年平均递增率为 \(r\)，则 \(r\ge0\)，且
\[
10(1+r)^{20}\le12
\]
即 \((1+r)^{20}\le1.2000\). 由于常用对数函数单调递增，
\[
20\lg(1+r)\le\lg1.2000=0.07918
\]
从而
\[
\lg(1+r)\le\frac{0.07918}{20}=0.003959
\]
因此严格的最大值在等号成立时取得：
\[
r=1.2^{1/20}-1\approx0.0091577558=0.91577558\%
\]
因此严格的最大递增率为 \(0.91577558\%\)；按题给对数表和通常保留百分数小数点后两位的报告精度，答案写作 \(0.92\%\).
\end{enumerate}
\end{solution}

\begin{problem}
\(ABCD-A_{1}B_{1}C_{1}D_{1}\) 为一正四棱柱，过 \(A\)，\(C\)，\(B_{1}\) 三点作一截面，求证：截面 \(ACB_{1}\perp\) 对角面 \(DBB_{1}D_{1}\).
\end{problem}

\begin{solution}
设 \(O=AC\cap BD\)，并记对角面 \(DBB_{1}D_{1}\) 为平面 \(\Pi\). 因为 \(ABCD\) 是正方形，所以两条对角线互相垂直，即
\[
AC\perp BD
\]
过 \(O\) 作直线 \(l\parallel BB_{1}\). 因为 \(O\) 在 \(DB\) 上，且平面 \(\Pi\) 包含直线 \(DB\) 和直线 \(BB_{1}\)，所以 \(l\subset\Pi\). 又因为正四棱柱的侧棱垂直于底面 \(ABCD\)，所以 \(l\perp AC\). 因此，直线 \(AC\) 垂直于平面 \(\Pi\) 内经过 \(O\) 的两条相交直线 \(BD\) 和 \(l\)，从而
\[
AC\perp\Pi
\]
而 \(AC\) 在截面平面 \(ACB_{1}\) 内，所以截面平面 \(ACB_{1}\) 与平面 \(DBB_{1}D_{1}\) 垂直，即
\[
ACB_{1}\perp DBB_{1}D_{1}
\]
\end{solution}

\begin{problem}
\begin{enumerate}
\item 设抛物线 \(y^{2}=4x\) 截直线 \(y=2x+k\) 所得的弦长为 \(3\sqrt{5}\)，求 \(k\) 的值.
\item 以本题（1）得到的弦为底边，以 \(x\) 轴上的点 \(P\) 为顶点做成三角形. 当这三角形的面积为 \(9\) 时，求 \(P\) 的坐标.
\end{enumerate}
\end{problem}

\begin{solution}
\begin{enumerate}
\item 将抛物线参数化为
\[
Q(t)=(t^2,2t)
\]
因为 \(y^2=4x\). 交点参数 \(t_1\)，\(t_2\) 满足直线方程，故
\[
2t=2t^2+k
\]
即
\[
t^2-t+\frac{k}{2}=0
\]
所以
\[
t_1+t_2=1
\]
两交点为 \(Q(t_1)\)、\(Q(t_2)\)，弦长的平方为
\[
\left|Q(t_1)Q(t_2)\right|^2=\left(t_1^2-t_2^2\right)^2+\left(2t_1-2t_2\right)^2=\left(t_1-t_2\right)^2\left(\left(t_1+t_2\right)^2+4\right)=5\left(t_1-t_2\right)^2
\]
由弦长为 \(3\sqrt5\)，得 \(5(t_1-t_2)^2=45\)，即 \((t_1-t_2)^2=9\). 又由韦达定理，
\[
(t_1-t_2)^2=(t_1+t_2)^2-4t_1t_2=1-4\cdot\frac{k}{2}=1-2k
\]
因此 \(1-2k=9\)，解得 \(k=-4\). 此时交点参数为 \(2,-1\)，确有两个实交点，且弦长为 \(3\sqrt5\)，检验成立.
\item 当 \(k=-4\) 时，弦所在直线为 \(y=2x-4\)，交点为
\(Q(2)=(4,4)\)，\(Q(-1)=(1,-2)\)
所以底边长为 \(3\sqrt5\). 设 \(P=(u,0)\)，则点 \(P\) 到底边所在直线 \(2x-y-4=0\) 的距离为
\[
h=\frac{|2u-4|}{\sqrt5}
\]
由三角形面积为 \(9\)，有
\[
9=\frac12\cdot3\sqrt5\cdot\frac{|2u-4|}{\sqrt5}=\frac32|2u-4|
\]
于是 \(|2u-4|=6\)，解得 \(u=5\) 或 \(u=-1\). 故
\[
P=(5,0)\quad\text{或}\quad P=(-1,0)
\]
\end{enumerate}
\end{solution}
